% GENERATED by reporting/make_tsae_report.py -- do not edit by hand.
% Source: metrics/outputs/{toy-temporal,gemma-2-2b*}/ -- frozen run results, not prose.
% Every caption is TODO on purpose: a caption is a claim about what to take away,
% and this project requires those to be written by a person.
% Regenerate:  python3 -m reporting.make_tsae_report
% Requires:    \input{thresholds_macros.tex} in the preamble (defines \MinJoint etc.),
%              and packages booktabs, graphicx, amsmath.


% ------------------------------------------------------------------------
% Toy model
% ------------------------------------------------------------------------

% [toy]
\begin{table}[tbp]
  \centering
  \small
  \caption{TODO(caption, human-written). The Temporal SAE on the small made-up world. Each row is one setting, run three times. The last column is the time rule. It has a value it cannot go above, $5.545$ here, which is what it shows when it has learned nothing at all. A row sitting at that value means the rule did nothing in that setting.}
  \label{tab:tsae-toy-conditions}
  \begin{tabular}{lcrrrrr}
    \toprule
    condition & $n$ & FVE & $L_0$ & parents in $B_0$ & InfoNCE (cos) & InfoNCE (inner) \\
    \midrule
    temporal, cosine & 3 & 0.983\,$\pm$\,0.004 & 1.80\,$\pm$\,0.29 & 2.67\,$\pm$\,0.47 & 5.217\,$\pm$\,0.046 & 5.435\,$\pm$\,0.016 \\
    temporal, inner product & 3 & 0.984\,$\pm$\,0.002 & 1.53\,$\pm$\,0.39 & 3.00\,$\pm$\,0.00 & 5.189\,$\pm$\,0.047 & 4.579\,$\pm$\,0.044 \\
    temporal, no term & 3 & 0.983\,$\pm$\,0.005 & 1.94\,$\pm$\,0.09 & 2.33\,$\pm$\,0.47 & 5.249\,$\pm$\,0.031 & 5.475\,$\pm$\,0.010 \\
    i.i.d., cosine & 3 & 0.986\,$\pm$\,0.003 & 2.25\,$\pm$\,0.32 & 2.67\,$\pm$\,0.47 & 5.575\,$\pm$\,0.002 & 5.546\,$\pm$\,0.001 \\
    i.i.d., inner product & 3 & 0.986\,$\pm$\,0.003 & 2.01\,$\pm$\,0.04 & 2.67\,$\pm$\,0.47 & 5.583\,$\pm$\,0.006 & 5.547\,$\pm$\,0.001 \\
    \bottomrule
  \end{tabular}
  \\[2pt] \footnotesize Each model is scored under both similarity functions; a model is comparable only to a baseline measured the same way. The tree has three parents.
\end{table}

% [toy-decorr]
\begin{table}[tbp]
  \centering
  \small
  \caption{TODO(caption, human-written). How many of the three parents end up in the first block, on two worlds. In the first world the parents are also the features that fire most often. In the second we made them rare instead. Everything else is the same, including how many features fire at once. The parents are learned in both worlds. Only the first block changes.}
  \label{tab:tsae-decorrelated}
  \begin{tabular}{lrrr}
    \toprule
    condition & parents, original & parents, decorrelated & distractors, decorrelated \\
    \midrule
    inner product & 3.00\,$\pm$\,0.00 & 0.67\,$\pm$\,0.47 & 3.33\,$\pm$\,0.47 \\
    cosine & 2.67\,$\pm$\,0.47 & 1.33\,$\pm$\,0.94 & 2.67\,$\pm$\,0.94 \\
    no term & 2.33\,$\pm$\,0.47 & 1.33\,$\pm$\,0.94 & 2.67\,$\pm$\,0.94 \\
    \bottomrule
  \end{tabular}
  \\[2pt] \footnotesize Out of three parents. All three are still learned somewhere in the dictionary in every run, so this is placement rather than a failure to learn them.
\end{table}

% ------------------------------------------------------------------------
% gemma-2-2b layer 12
% ------------------------------------------------------------------------

% [gemma]
\begin{table}[tbp]
  \centering
  \small
  \caption{TODO(caption, human-written). Two SAEs at layer 12 of \texttt{gemma-2-2b}, measured the same way on the same text. The blocks differ a lot in size, so the number of links found cannot be compared directly. The density column can: it is the share of all possible pairs that became a link. The shape column gives the block sizes behind each row.}
  \label{tab:tsae-gemma-pipeline}
  \begin{tabular}{llrrrrr}
    \toprule
    block pair & shape & edges & density & superparents & joint-child cov. & dropped by $\MinJoint$ \\
    \midrule
    Temporal SAE 0$\rightarrow$1 & 3276$\times$13108 & 1,621 & 3.77e-05 & 0 & 0.381 & 2,281 \\
    Matryoshka 0$\rightarrow$1 & 128$\times$384 & 1,473 & 3.00e-02 & 2 & 0.506 & 2 \\
    Matryoshka 1$\rightarrow$2 & 384$\times$1536 & 621 & 1.05e-03 & 0 & 0.223 & 93 \\
    \bottomrule
  \end{tabular}
  \\[2pt] \footnotesize Raw counts do not compare across pairs of different shape; density does. The Matryoshka $1\rightarrow2$ row is the closest in shape, not a matched control.
\end{table}

% [gemma-sres]
\begin{table}[tbp]
  \centering
  \small
  \caption{TODO(caption, human-written). Our strictest test, at layer 12 of \texttt{gemma-2-2b}. It does not ask whether the two features fire together. It asks whether the parent tells us something about the child that the child does not already tell us by itself. The last column is the share of tested links that passed.}
  \label{tab:tsae-gemma-sres}
  \begin{tabular}{lrrr}
    \toprule
    block pair & edges scored & pass & fraction \\
    \midrule
    Temporal SAE $0\rightarrow1$ & 1,318 & 770 & 0.584 \\
    Matryoshka $0\rightarrow1$ & 1,111 & 22 & 0.020 \\
    Matryoshka $1\rightarrow2$ & 613 & 65 & 0.106 \\
    Matryoshka $2\rightarrow3$ & 1,720 & 212 & 0.123 \\
    \bottomrule
  \end{tabular}
  \\[2pt] \footnotesize No child was untestable in either run, so neither fraction is thinned by missing negatives.
\end{table}

% [thresholds]
\begin{table}[tbp]
  \centering
  \small
  \caption{TODO(caption, human-written). What happens when we raise our main cut-off, at layer 12 of \texttt{gemma-2-2b}. Raising a cut-off should keep only the stronger links. The column on the right is the share of links that could be explained by the two features simply being common. For Matryoshka that share goes up, which is the opposite of what a stricter cut-off should do.}
  \label{tab:tsae-threshold-reversal}
  \begin{tabular}{lrrrr}
    \toprule
    $\EdgeTau$ & MSAE edges & chance share & T-SAE edges & chance share \\
    \midrule
    0.1 & 1,877 & 0.65 & 13,347 & 0.02 \\
    0.5 $\leftarrow$ operating & 519 & 0.77 & 2,670 & 0.00 \\
    0.6 & 375 & 0.85 & 1,620 & 0.00 \\
    0.7 & 314 & 0.94 & 1,019 & 0.00 \\
    0.9 & 281 & 1.00 & 432 & 0.00 \\
    \bottomrule
  \end{tabular}
  \\[2pt] \footnotesize Chance share is the fraction of accepted edges with PMI below 0.5. The two columns differ in first-block width, 128 against 3{,}276, which is confounded with the architecture here.
\end{table}
